\section{Linearized measures and signed excess}\label{sec:linearized}

The centered blow-up of Proposition~\ref{prop:centered-blowup} has small
height, small tilt, and negligible mean curvature at the height scale.
We now retain its second-order information as measures on the limiting
plane. This formulation allows tilt to concentrate: no compactness theorem
for a limiting collection of graphs is required.
A related use of measure-valued limits for stationary two-valued graphs
appears in the generalized-gradient Young measures and linear measure
solutions of Hirsch and Spolaor \cite[Sections~2--3]{HSgraphs}.
The measure equations needed here are derived directly below.

For a concrete model, let $u:\Omega\to\R^n$ be a smooth harmonic
map on an open subset of $\R^m$. Its squared height, height flux,
and gradient energy define the measures
\[
 \begin{aligned}
 \nu_u&=|u|^2\,\dd y,
 & (f_u)_a&=u\cdot\partial_a u\,\dd y,\\
 (\beta_u)_{ab}&=\partial_a u\cdot\partial_b u\,\dd y,
 & \mathfrak m_u&=\tfrac12|Du|^2\,\dd y,
 \end{aligned}
\]
where $1\le a,b\le m$. Their block matrix is a Gram matrix, and
harmonicity gives
$\nabla\nu_u=2f_u$, $\diver f_u=\tr\beta_u$, and
$\diver\beta_u=\nabla\mathfrak m_u$ in distributions.
Moreover $2\mathfrak m_u=\tr\beta_u$.
The proposition below retains these identities and block positivity,
with the last equality replaced by a one-sided inequality. The model
explains the roles of the measures; the actual limit is not assumed
to admit such a harmonic-map representation.

\begin{proposition}[The limiting measure equations]\label{prop:linearized}
Let $V_i$ be stationary integral $m$-varifolds converging locally to
$Q|\Pi|$, where $Q$ is a positive integer and $\Pi$ is an $m$-plane.
Let $S_i$ be reference graphs converging to $\Pi$ in every fixed finite
derivative order used below, on a normal cylinder with fixed interior
room over a connected disk $\Omega\subset\Pi$. The cylinders contain all
support used below and have a support-free gap at their artificial vertical
boundaries above each compact base subset. Let $\mu_i$ denote the weight
of $V_i$ restricted to this cylinder, $\pi_i$ nearest projection onto
$S_i$, and $z_i=X-\pi_iX$. All pushforwards below use this restriction.
Use the horizontal and vertical tangent blocks relative to $S_i$ to set
$B_i=I-(P_T)_{hh}$, $e_i=\tr B_i$,
$D_i=(\pi_i)_\#(e_i\mu_i)$, and
$M_i=(\pi_i)_\#\mu_i-Q\vol_{S_i}$.
Suppose $a_i\downarrow0$ and, on each compact subdisk with interior room,
\begin{gather*}
 |M_i|+D_i=O(a_i^2),\qquad |z_i|\le Ca_i\quad\hbox{on the support},\\
 |\mathcal H_{S_i}|+|\nabla^\perp\mathcal H_{S_i}|=o(a_i),
 \qquad a_i^{-1}(\pi_i)_\#(z_i\mu_i)\rightharpoonup0.
\end{gather*}
Assume also that $a_i^{-2}(\pi_i)_\#(|z_i|^2\mu_i)$ has a positive
uniform lower mass bound on a fixed compact subset of $\Omega$.
Identify the bases by their converging graph coordinates. After passage
to a subsequence there are Radon measures satisfying
\begin{equation}\label{eq:limit-measures}
\begin{aligned}
 a_i^{-2}(\pi_i)_\#(|z_i|^2\mu_i)&\rightharpoonup\nu,&
 a_i^{-2}(\pi_i)_\#(B_i\mu_i)&\rightharpoonup\beta,\\
 a_i^{-2}D_i&\rightharpoonup d_0=\tr\beta,&
 a_i^{-2}M_i&\rightharpoonup\mathfrak m,\\
 a_i^{-2}(\pi_i)_\#\bigl(z_i\cdot(P_T)_{vh}\,\mu_i\bigr)
   &\rightharpoonup f.
\end{aligned}
\end{equation}
Here $\nu\ne0$ and $d_0$ are nonnegative measures, $f$ is vector-valued,
$\beta$ is symmetric matrix-valued, and $\mathfrak m$ is signed. They obey
\begin{equation}\label{eq:limit-psd}
 \begin{pmatrix}\nu&f^{\mathsf t}\\ f&\beta\end{pmatrix}\ge0
\end{equation}
as a matrix-valued measure, and
\begin{equation}\label{eq:limit-identities}
 \nabla\nu=2f,\qquad \diver f=d_0,\qquad
 \int_\Omega\diver X\,\dd\mathfrak m
   =\int_\Omega DX:\dd\beta
 \quad (X\in C_c^\infty(\Omega;\R^m)).
\end{equation}
Finally, locally on $\Omega$,
\begin{equation}\label{eq:limit-sign}
             2\mathfrak m\le d_0,\qquad
             |\mathfrak m|\le C d_0.
\end{equation}
The data supplied by Proposition~\ref{prop:centered-blowup} satisfy these
hypotheses, with a fixed disk containing the closed unit ball of density
centers.
\end{proposition}

\begin{proof}
All estimates are on compact subdisks with a fixed enlargement. Graph
volume factors are included when measures are transported to $\Omega$,
and tensor components are taken in smoothly converging orthonormal frames.
The assumed mass bounds give weak compactness. The projection identity
\[
 (P_T)_{hv}(P_T)_{vh}=B_i-B_i^2\le B_i
\]
shows both that the flux measures have bounded mass and that the block
matrix before passage to the limit is positive semidefinite. This proves
\eqref{eq:limit-psd}; the lower second-moment bound gives $\nu\ne0$.

\subsection*{Vanishing geometric errors}

We record the horizontal first variation carefully, since the height
normalization may be much smaller than the curvature of the reference
graph. Suppress $i$ temporarily. For a tangent field $Y$ on $S$, let
$J_z=I-\mathcal S_z$, where $\mathcal S_z$ is the shape operator, and
lift $Y$ to the ambient field $J_zY$. The curved first-variation calculation
in \cite[Lemma~2.2]{AE} gives
\begin{equation}\label{eq:horizontal-error}
 P_T:D(J_zY)=\diver_S Y-B:\nabla^S Y+\mathcal E_Y,
\end{equation}
with the following sharper error in the present setting. Choose
$\varepsilon_i\to0$ dominating the fixed geometric derivatives of $S_i$
that tend to zero and such that
$|\mathcal H_{S_i}|+|\nabla^\perp\mathcal H_{S_i}|\le\varepsilon_i a_i$.
Then
\begin{equation}\label{eq:horizontal-error-bound}
 |\mathcal E_Y|\le C\varepsilon_i\left\{
 |z_i|e_i|\nabla^{S_i}Y|
 +(a_i|z_i|+|z_i|^2+\sqrt{e_i}\,|z_i|)|Y|\right\}.
\end{equation}
Indeed, in frames whose tangent and normal connections vanish at the
point of calculation, the horizontal derivative has blocks
\[
 D(J_zY)_{hh}=J_z(\nabla^SY)J_z^{-1}-\mathcal K_zJ_z^{-1},
 \qquad \mathcal K_z(a)=(\nabla_a\mathcal S)_zY.
\]
Conjugation preserves the trace of $\nabla^SY$, and its pairing with $B$
has error $O(\varepsilon_i|z|e|\nabla^SY|)$. Codazzi gives
$\tr\mathcal K_z=\langle\nabla_Y^\perp\mathcal H_S,z\rangle$.
The remaining horizontal terms are bounded by
$C\varepsilon_i(a_i|z|+|z|^2+e|z|)|Y|$.
The symmetric sum of the off-diagonal blocks is, when paired with a
tangent vector $a$ and a normal vector $v$,
\[
 \big\langle a,J_z^{-1}[\mathcal S_z,\mathcal S_v]Y\big\rangle.
\]
It is $O(\varepsilon_i^2|z||a||v||Y|)$; pairing with $(P_T)_{vh}$
costs at most $C\sqrt e$. This proves
\eqref{eq:horizontal-error-bound}, including in arbitrary codimension.

The corresponding normal formula, obtained from $D\pi=J_z^{-1}P_h$, is
\begin{equation}\label{eq:normal-trace}
 P_T:Dz=e-\mathcal H_S\cdot z
       +O(|\mathrm{II}_S|^2|z|^2+|\mathrm{II}_S|e|z|).
\end{equation}
This is the normal trace computation used in the minimal-fitting
argument of \cite[Equation~(7.9)]{AE}. At normalization $a_i^{-2}$, every
error in \eqref{eq:horizontal-error-bound} and
\eqref{eq:normal-trace} tends to zero in mass: use $|z_i|\le Ca_i$,
$\int e_i\,\dd\mu_i=O(a_i^2)$, and Cauchy--Schwarz for
$\int |z_i|\sqrt{e_i}\,\dd\mu_i$.

Stationarity tested with $J_{z_i}Y$ now yields the last identity in
\eqref{eq:limit-identities}; the reference-volume integral of a divergence
is zero. Testing instead with $|z_i|^2J_{z_i}Y$ gives
$\nabla\nu=2f$. Terms containing both $B_i$ and $|z_i|^2$ disappear,
and the differentiated height contributes
$2z_i\cdot(P_T)_{vh}J_{z_i}Y$; replacing $J_{z_i}$ by the identity has
vanishing normalized error. Finally, the normal test $\varphi z_i$,
with $\varphi$ a smooth base function, gives
$\int\varphi\,\dd d_0+\int\nabla\varphi\cdot\dd f=0$ by
\eqref{eq:normal-trace}. Thus all equations in
\eqref{eq:limit-identities} hold.

\subsection*{Applying signed excess at the height normalization}

It remains to prove \eqref{eq:limit-sign}. There are two requirements
for applying Theorem~\ref{thm:signed-input} at this normalization.
First, its geometric error $\Lambda_i$ must be $o(a_i^2)$; in target
units its height-square term is only $O(a_i^2)$. We achieve this by
changing physical units. Second, the physical chart radius and support
height must be small relative to the local excess budget. We meet that
condition by dividing the base into cells whose radii are powers of
$a_i$. These choices do not require a rate relating $a_i$ to the
original target scales.

Fix a compact test support and an interior enlargement. Dilate physical
space by $d_i'\downarrow0$ sufficiently slowly that all required finite
geometric derivatives remain bounded and
\[
 \frac{|\mathcal H_{S_i}|}{d_i'a_i}\longrightarrow0,
 \qquad
 \frac{|\nabla^\perp\mathcal H_{S_i}|}{(d_i')^2a_i}\longrightarrow0.
\]
Such a choice exists because each numerator divided by $a_i$ tends to
zero, and all the finitely many graph derivatives tend to their flat
values. To obtain a uniform normal tube after dilation, write the graph
as $g_i$ and extend it using a fixed smooth cutoff $\chi$ equal to one
on the enlargement. The new physical graph and its derivatives are
\[
 g_i^{\rm phys}(v)=d_i'(\chi g_i)(v/d_i'),\qquad
 D^j g_i^{\rm phys}(v)
   =(d_i')^{1-j}D^j(\chi g_i)(v/d_i').
\]
The cutoff graph tends to zero through all the required finite orders,
including order zero. Choosing $d_i'$ still more slowly keeps these
derivatives bounded and the second derivatives small, which provides
the fixed tube. The mean-curvature bounds are needed only on the
unchanged inner region. The cylinder restriction has a support gap at
its artificial vertical boundary, so this operation introduces no
first-variation term.

In the new physical units, the geometric measure $\Lambda_i$ charges
\[
 (d_i')^2|z_i|^2+(d_i')^{-2}|\mathcal H_{S_i}|^2
                  +(d_i')^{-4}|\nabla^\perp\mathcal H_{S_i}|^2.
\]
Expressed with mass returned to target units, it therefore satisfies
$\Lambda_i=o(a_i^2)$. The measures $M_i,D_i$ in target normalization
are unchanged. Thus the geometric error will disappear after division
by $a_i^2$; we now arrange the local hypotheses for signed excess.

Use $\zeta$ and $\kappa$ fixed in Equation~\eqref{eq:parameter-order},
and cover the test support by cells of
target radius $\lambda_i=a_i^\zeta$, with bounded overlap of their fixed
enlargements. In particular $3\kappa<\zeta$. On each enlarged cell use a normalized budget
$b_{i,j}$ for $|M_i|+D_i+\Lambda_i$, enlarged by fixed constants and
with floor $a_i^3$. Thus, uniformly in the cells,
\begin{equation}\label{eq:signed-cell-budget}
 a_i^3\le b_{i,j}\le C a_i^2\lambda_i^{-m}+a_i^3=o(1),
 \qquad \sum_j\lambda_i^m b_{i,j}=O(a_i^2).
\end{equation}
Fresh normal charts on these cells have physical radius
$O(d_i'\lambda_i)$ and physical support height $O(d_i'a_i)$.
Since $3\kappa<\zeta<1$, these are at most $b_{i,j}^{\kappa}$ for
all large $i$, with the fixed factors absorbed by strictness of the
inequality. All physical hypotheses of signed excess are satisfied.

Partition a nonnegative smooth test by a smooth partition of unity at
scale $\lambda_i$. After the fixed normalization of its size and
Lipschitz constant, the signed-excess error on a cell is
$o(1)b_{i,j}$, uniformly because $b_{i,j}\to0$ uniformly. Multiplying
by $\lambda_i^m$, summing \eqref{eq:signed-cell-budget}, and dividing
by $a_i^2$ proves $2\mathfrak m\le d_0$.

For the absolute bound, apply Lemma~\ref{lem:scalar-input}, from
\cite[Lemma~2.5]{AE}, in the same physical units on any fixed interior
ball and its enlargement. Smooth mass averages identify its integer as
$Q$ for all large $i$. Since $\Lambda_i=o(a_i^2)$, passage to the
limit gives, after harmless fixed changes of radii,
\[
       |\mathfrak m|(B_u(y))\le C d_0(B_{3u}(y)).
\]
This inflated-ball estimate implies domination as measures, without a
doubling assumption. For a compact $K\subset O\Subset\Omega$, cover
$K$ by sufficiently fine equal-radius balls whose triple enlargements
lie in $O$ and have bounded overlap. Summing gives
$|\mathfrak m|(K)\le C d_0(O)$. Outer regularity as $O$ decreases to
$K$, followed by inner regularity, proves $|\mathfrak m|\le C d_0$.
This completes the proof.
\end{proof}
